\subsection{The complete global central localization map and its joint quotient}

\begin{statement}{E.93}{Proved}{Every central localization fibre and its exact action on global class-two sections}
Keep the cubic curve, rational tangent, genuine arithmetic action,
ten-generator lattice $V$, and class-two group $N_2$ of E.37 and
E.74. Put $M=\Gamma_2N/\Gamma_3N=\bigwedge^2V$, with the
actual Hall convention of E.74, and write
\[
 \begin{gathered}
 \mathcal P_2=M/(I-T_2)M^{J_2}
       \simeq\mathbb Z_3^{24}\oplus C_3^9\oplus C_9^5,\\
 \mathcal K_2=\ker\bigl((I-J_2)(T_2-7I):M\longrightarrow M\bigr)
       \simeq\mathbb Z_3^{22}.
 \end{gathered}
 \tag{E93.1}
\]
All cohomology here is continuous integral cohomology. For a set
of rational primes $\mathcal S$, let
\[
 \mathcal R_{\mathcal S}(\mathbb Z_3)
   =\{(a_p)_{p\in\mathcal S}:\text{for every }n,
             \ a_p\in3^n\mathbb Z_3\text{ for all but finitely many }p\}
 \tag{E93.2}
\]
carry the inverse-limit topology of the discrete groups
$\bigoplus_{p\in\mathcal S}\mathbb Z/3^n$.

Let $w$ be the primitive integral Tate vector of E.37:
$w=(3f_1+U_1)\wedge(3f_2+U_2)$. Its actual based actions
are $T_2w=7w$ and $J_2w=-w$. Write $[w]$ for its class in
$\mathcal P_2$; no choice of a free basis vector for that
class is made. In the actual global splitting of E.37 and
the actual local coordinates of E.75 the restriction map is
\[
 \begin{gathered}
 H^1(\mathbb Q,M)
       =\mathcal P_2\oplus\mathcal R_{\{p\ne3\}}(\mathbb Z_3),\\
 H^1(\mathbb Q_3,M)=\mathcal P_2\oplus\mathcal K_2,\\
 \operatorname{res}\bigl(z,(a_p)\bigr)
       =\bigl(z+b\mathcal L(a)[w],\ \alpha\mathcal L(a)w\bigr),\\
 \mathcal L(a)=\sum_{p\ne3}a_p\ell_p,
 \qquad \ell_p=\frac{\log\langle p\rangle}{\log4}\in\mathbb Z_3,
 \qquad \ell_2=\tfrac12.
 \end{gathered}
 \tag{E93.3}
\]
The brackets $\langle p\rangle$ remove the sign of the
three-adic unit. The constants $b\in\mathbb Z_3$ and
$\alpha\in\mathbb Z_3^\times$ are the actual unchanged
marked Kummer evaluations of E.90:
$\kappa(3)=(1,0)$ and $\kappa(4)=(b,\alpha)$.
Their numerical values are not assigned.

The entire image, kernel and cokernel of this specific
continuous homomorphism are
\[
 \begin{aligned}
 \operatorname{im}(\operatorname{res})
    &=\mathcal P_2\oplus\mathbb Z_3w
      \simeq\mathbb Z_3^{25}\oplus C_3^9\oplus C_9^5,\\
 \ker(\operatorname{res})
    &\simeq\mathcal R_{\{p\ne2,3\}}(\mathbb Z_3),\\
 \operatorname{coker}(\operatorname{res})
    &=\mathcal K_2/\mathbb Z_3w\simeq\mathbb Z_3^{21}.
 \end{aligned}
 \tag{E93.4}
\]
The isomorphism for the kernel sends an arbitrary restricted
family $(a_p)_{p\ne2,3}$ to
\[
 z=0,\qquad
 a_2=-2\sum_{p\ne2,3}a_p\ell_p,
 \tag{E93.5}
\]
with the other coordinates unchanged. A continuous additive
section onto the image sends $(u,t w)$ to
\[
 z=u-\frac{bt}{\alpha}[w],\qquad
 a_2=\frac{2t}{\alpha},\qquad a_p=0\quad(p\ne2,3).
 \tag{E93.6}
\]
Adding (E93.5) to (E93.6) gives every localization fibre.

Here are exact coordinates for the cokernel. Let $K$ and
$K^{\ell}$ be E.75's retained $45\times22$ saturated kernel
basis and its left inverse. The actual column $k=K^{\ell}w$ is
\[
 \begin{split}
 k={}&(-2,-4,0,-4,-2,0,0,1,-2,3,0,\\
     &\hspace{19mm}2,4,6,6,-2,-2,1,-4,-1,2,-3)^t.
 \end{split}
 \tag{E93.7}
\]
For $a=K\eta$, define
\[
 t=-\eta_1/2,\qquad
 \xi_j=\eta_j+(k_j/2)\eta_1\quad(2\le j\le22).
 \tag{E93.8}
\]
Then $a=t w+\sum_{j=2}^{22}\xi_j K e_j$.
Thus $(\xi_2,\ldots,\xi_{22})$ is the full quotient
map and $\xi\mapsto\sum_{j=2}^{22}\xi_j K e_j$ is its
continuous integral section. These maps, including their
full $45$-coordinate matrices, are retained in the certificate.

For the two actual quotient sections used in E.85, let
$F_2:M\to A=\mathbb Z_3^2$ be its exact alternating
coefficient map. Give $H^1(\mathbb Q_3,A)$ the complete
coordinates $(\rho,\sigma,\epsilon)$ of E.85, including
geometric conjugacy. The composite map is exactly
\[
 \begin{gathered}
 H^1(\mathbb Q,M)\xrightarrow{\operatorname{res}}
 H^1(\mathbb Q_3,M)\xrightarrow{F_{2,*}}
 H^1(\mathbb Q_3,A),\\
 \bigl(z,(a_p)\bigr)\longmapsto(\rho(z),0,0),
 \qquad \rho(z)=(F_{2,1}-F_{2,2})z,\\
 \operatorname{im}=\mathbb Z_3\oplus0\oplus0,\\
 \ker=\mathcal R_{\{p\ne3\}}(\mathbb Z_3)
       \oplus\mathbb Z_3^{23}\oplus C_3^9\oplus C_9^5.
 \end{gathered}
 \tag{E93.9}
\]
The actual horizontal splitting vector is
$z_*=-\frac12[f_2,f_1]$, with
$F_2z_*=(1/2,-1/2)$. In the Smith coordinates
$p=L_Gz$ of E.85, retain all fourteen finite coordinates
and the final twenty-three free coordinates. The missing
free coordinate is obtained from
\[
 \begin{gathered}
 p_{22}=\sum_{j=2}^{24}r_jp_{21+j}-\rho,\\
 r=(-1,-1,-1,0,0,0,0,-3,3,0,-1,0,
       -1,0,0,3,-3,0,0,-2,-3,-3,-1,0).
 \end{gathered}
 \tag{E93.10}
\]
The exact full inverse matrices occur in the certificate.

These maps apply to the actual nonabelian fibres as follows.
Let $\mathcal I\subset H^1(\mathbb Q,V)$ be the full image
of $H^1(\mathbb Q,N_2)$ computed by the explicit obstruction
criterion in E.37. For each $\gamma\in\mathcal I$, let
$\mathfrak F_\gamma$ be its nonempty global section fibre.
Its restriction to the local intrinsic fibre is equivariant
for (E93.3). After choosing any actual class in
$\mathfrak F_\gamma$ as origin, the image is its translate
of the group in (E93.4), and every fibre is the group in
(E93.5). The image of $\mathfrak F_\gamma$ in the pair of
local quotient-section classes is its origin's translate
of $\mathbb Z_3\oplus0\oplus0$, and every fibre has exactly
the kernel in (E93.9). A change of global origin translates
these coordinates by its unique class in $H^1(\mathbb Q,M)$.
No arbitrary origin is assigned a numerical joint value.

This computes central variations on every actual global
class-two fibre. Evaluating the particular joint origin for
every global abelian input, and all higher nonabelian lifts,
remain the stated steps of E.23. The manuscript passages
are section51.12, folios717--718, scans136--138; section51.15,
folios725--727, scans156--160; and section51.16, folio734v,
scan174, with the exact comparisons of E.37 and E.85.
\end{statement}
\begin{proof}
\emph{The genuine coefficient line and its basis transport.}
The module $M$ is the central kernel in
$1\to M\to N_2\to V\to1$, with conjugation action induced
by the retained rational tangential section. E.37 constructs
the integral vector $w$ in the lexicographic exterior basis
$(q_i\wedge q_j)_{i<j}$. E.74 uses Hall generators
$[q_a,q_b]$ with $a>b$. We use the explicit signed permutation
matrix $S$ whose row for $[q_a,q_b]$ has entry $-1$ in
column $(b,a)$ and zero elsewhere. Thus
\[
 SS^t=I,\qquad
 T_2S=S\bigwedge^2T_1,\qquad J_2S=S\bigwedge^2J_1.
 \tag{E93.11}
\]
These are equalities of maps between the two specified
free $\mathbb Z_3$-modules. The certificate checks every
entry and transports $w$ by $S$ before applying any Hall
matrix. Its Hall coordinate numbered nine is $-1$, proving
primitivity in $M$ directly. The actual arithmetic action
then gives $T_2w=7w$ and $J_2w=-w$. In particular the map
$\mathbb Z_3(1)\to M$, $1\mapsto w$, is equivariant for
the full global action, as proved in E.37.

\emph{The global and local decompositions and restriction.}
E.37.20 gives an exact continuous localization sequence
whose residue module is (E93.2) for $p\ne3$. Its proof uses
the actual inertia maps: the $p$-eigenlattice of Frobenius
in $M$ is exactly $\mathbb Z_3w$. The explicit splitting is
\[
 (a_p)\longmapsto\sum_{p\ne3}a_p\kappa(p)w.
 \tag{E93.12}
\]
Every reduction modulo $3^n$ in this sum has only finitely
many nonzero terms. These reductions are compatible, so the
sum is a uniformly convergent continuous cocycle. It has
the specified inertia residues at all primes. Its complement
is the unramified term evaluated at the actual free
generator $\tau$, namely $\mathcal P_2$ in (E93.1).
The image of this unramified cocycle under the actual local
quotient map of E.74 is $(z,0)$: the marked relative
generator dies and the marked $\tau$ is unchanged. This
identifies the first summand through an actual cocycle map.

E.75 evaluates the local normalized cocycles as
$\mathcal P_2\oplus\mathcal K_2$, with horizontal value
modulo $(I-T_2)M^{J_2}$ and literal relative value in
$\mathcal K_2$. The evaluated Kummer comparison of E.90
identifies $H^1(\mathbb Q_3,\mathbb Z_3(1))$ with
$\mathbb Z_3\kappa(3)\oplus\mathbb Z_3\kappa(4)$.
For the rational prime $p\ne3$ its valuation coordinate
is zero and its unit coordinate is $\ell_p$. Hence
\[
 \operatorname{res}\kappa(p)
       =\ell_p\kappa(4)
       =(b\ell_p,\alpha\ell_p)
 \tag{E93.13}
\]
in the actual marked evaluations. This is an equality of
cohomology classes with the same chosen generators.
Multiplying by the equivariant vector $w$, then adding the
unramified class, proves (E93.3).

For completeness, $\log4$ has three-adic valuation one,
and $\log\langle p\rangle\in3\mathbb Z_3$; therefore
every $\ell_p$ is integral. The identity $(-2)^2=4$ in
$1+3\mathbb Z_3$ gives $\ell_2=1/2$. The uniform residue
condition (E93.2) makes $\mathcal L$ converge and makes it
continuous. Thus $a_2=2t$ with all other entries zero is a
continuous additive section of $\mathcal L$ onto
$\mathbb Z_3$.

The second coordinate of (E93.3) now has image exactly
$\mathbb Z_3w$, since $\alpha$ is a unit. Its first
coordinate can be assigned freely by changing $z$.
This proves the image assertion and verifies (E93.6)
by literal substitution. In the kernel, the torsion-free
line $\mathbb Z_3w$ forces $\mathcal L(a)=0$ and then
$z=0$. Solving its single unit equation using
$\ell_2=1/2$ gives precisely (E93.5), including its
inverse projection and topology.

\emph{The full integral quotient map.}
The matrices $K,K^{\ell}$ satisfy $K^{\ell}K=I$;
E.75 proves that $K$ maps $\mathbb Z_3^{22}$ onto
the entire kernel $\mathcal K_2$. Their exact product
with the transported $w$ is the column (E93.7), and
$Kk=w$. Replace the first column of the identity
matrix of order22 by $k$. Its determinant is $-2$,
a three-adic unit, so its matrix and inverse are
integral over $\mathbb Z_{(3)}$. The inverse equations
are precisely (E93.8). They prove directly that the
relative Tate line is a direct summand and give the
entire rank21 cokernel in (E93.4).

This relative primitivity does not enter as an
unproved primitivity assertion for the horizontal
class $[w]$. Indeed the program evaluates $L_Gw$;
its free entries are all divisible by three and
at least one is not divisible by nine. Formula
(E93.6) uses the actual class $[w]$ with its entire
Smith vector, so its correctness retains this
divisibility and all fourteen finite coordinates.

\emph{The exact joint coefficient map.}
E.85 constructs $F_2$ from the four actual unit
characters, with image
$L_2=\{(x,y):x+y\in3\mathbb Z_3\}\subset A$.
It proves the genuine arithmetic intertwiners
$F_2T_2=49F_2$ and $F_2J_2=\operatorname{swap}F_2$.
Applying the first equality to $w$ gives
\[
             7F_2w=49F_2w,\qquad F_2w=0,
 \tag{E93.14}
\]
the last implication following from the torsion-free
target and the nonzero integer42. The retained
matrix multiplication verifies the same zero
column after the exact transport (E93.11).

The full map on local cohomology in E85.12 is
\[
 ([z],a)\longmapsto
       \bigl((F_{2,1}-F_{2,2})z,\ F_{2,1}a,\ 0\bigr).
 \tag{E93.15}
\]
Its last entry is zero because $F_2z\in L_2$;
the actual central conjugator removes its mean
modulo three as in E85.25. Substituting (E93.3)
and (E93.14) proves (E93.9). At the cocycle level
every summand $a_p\kappa(p)w$ already has zero
image under $F_2$; no limiting ambiguity is used.

For every Smith coordinate the exact identity
$L_G(I-T_2)(I+J_2)R_G
  =\operatorname{diag}(1^7,3^9,9^5,0^{24})$
and both retained inverse matrices are checked.
Multiplying $\rho$ by $L_G^{-1}$ gives zero in
the first twenty-one columns and the row $r$
in (E93.10) in the last twenty-four. Its first
coefficient is $-1$. Solving that equation proves
the full $23$-dimensional free kernel while leaving
all $C_3^9\oplus C_9^5$ coordinates arbitrary.
The vector $z_*$ gives the displayed continuous
additive section of $\rho$. Combining these
coordinates with every restricted prime family
proves the entire kernel and every inverse fibre
in (E93.9).

\emph{Transport to actual section classes.}
The class-two central extension identifies the
difference between any two lifts of the same
abelian cocycle with a unique $M$-cocycle modulo
central coboundaries. Changing the representative
of the fixed abelian class introduces its
stabilizer $H^0(\mathbb Q,V)$. This group is zero
by E.36. The corresponding local stabilizer is
zero by E.74, and the target abelian stabilizer
is zero because its actual $\tau$-operator is
$7I$, as verified in E.85. Thus no extra
stabilizer quotient changes any of the displayed
coefficient fibres. The same cocycle multiplication
commutes with restriction and with the two
quotient homomorphisms, proving equivariance
and all the stated torsor assertions.

The choice of an origin is a translation within
the already nonempty fibre over $\gamma\in\mathcal I$.
Its change is the unique central cohomology
class just constructed, whose restriction and
joint image are (E93.3) and (E93.9). This gives
the exact ambiguity as well as both transfer maps.

All finite matrices in this proof are over
$\mathbb Z_{(3)}$ with the displayed integral
inverses. The restricted sums are continuous in
the topology (E93.2). The certificate checks the
two basis transports and all matrix identities,
the symbolic inverse retaining $b,\alpha$, $1134$
finite-precision localization identities, and
$1458$ full relative and horizontal coordinate
round trips. The algebraic identities and uniform
convergence above prove the full integral statements.
\end{proof}
